\par\Needspace{7\baselineskip}\bigskip
\begin{center}
\pdfbookmark[2]{Gauss an Olbers, Göttingen, 8. April 1813}{pallas-brief-18130408}
\textsc{Gauss} an \textsc{Olbers}, Göttingen, 8. April 1813.
\end{center}
\SourceLeaf{428}\noindent\hspace*{\parindent}. . . . .\quad Mit den Pallasstörungen bin ich ziemlich vorgerückt; es fehlt jetzt nur noch ein Element, die Epoche. Beim Perihel sind (vom {\AstroSymbols ♃} herrührend) 183 Gleichungen, die über \(0''{,}1\) betragen. Wenn die Epoche auch beseitigt ist, unternehme ich vielleicht noch wenigstens eine \textls[180]{Probe} von Tafeln; 800 Gleichungen schrecken freilich ab, aber der Umstand \(7\text{\AstroSymbols ⚴}=18\text{\AstroSymbols ♃}\) (quam proxime) erlaubt eine grosse Abkürzung. Inzwischen beendigt Hr. \textsc{Nicolai} wohl die Störungen durch {\AstroSymbols ♄}, und vielleicht unternimmt ein anderer sehr geschickter Schüler von mir, Hr. \textsc{Encke}, die Störungen durch {\AstroSymbols ♂}. Erlaubt es Ihre Zeit, so helfen Sie doch auch nächstens die {\AstroSymbols ⚴} aufsuchen.\quad . . . . .
\par\Needspace{7\baselineskip}\bigskip
\begin{center}
\pdfbookmark[2]{Gauss an Olbers, Göttingen, 2. Julius 1813}{pallas-brief-18130702}
\textsc{Gauss} an \textsc{Olbers}, Göttingen, 2. Julius 1813.
\end{center}
\noindent\hspace*{\parindent}. . . . .\quad Die Störungen der Pallas durch Jupiter, so weit ich sie zu berechnen die Absicht hatte, sind jetzt zum grössten Theil vollendet. Von den numerischen Rechnungen für die bisherigen Oppositionen ist aber noch viel zurück.\quad . . . . .
\par
